\section{Restliche Ideen}
\begin{frame}{Und jetzt?}
    \begin{itemize}
        \item Sind aber an $\left|\E[h(\mathbf{Q}(\X))]-\E[h(Z)]\right|$ interessiert. 
        \item $h = \indicator{(-\infty, x_1] \times \cdots \times (-\infty, x_m]}$ natürlich \textbf{nicht} differenzierbar. 
        \pause
        \item Wir machen es also differenzierbar!
        \item Mit Hilfe eines \textbf{gaußschen} Glättungsoperators!
        \[
        h_t(x) = \int_{\R^m} h(\sqrt{t}y + \sqrt{1-t}x) \mathcal{N}(0, \operatorname{Id})(dy), \quad t \in (0,1), 
        \]
        \pause
        \item \textbf{Fakt:} Sei $\H(\R^m)$ Menge der Indikatoren konvexer Mengen auf $\R^m$. Ist $Z$ \textit{standardnormal}, so gilt für jedes Wahrscheinlichkeitsmaß $Q$ auf $\R^m$ gilt für jedes $W \sim Q$ bereits \begin{align*}
        &\sup_{h \in \H(\R^m)} \left|\E[h(W)] - \E[h(Z)]\right| \\ & \qquad \qquad \qquad \leq \frac{4}{3} \sup_{h \in \H(\R^m)} \left(\left|\E[h_t(W)] - \E[h_t(Z)]\right| + 2 \sqrt{m}\sqrt{t}\right)
        \end{align*}
    \end{itemize}
\end{frame}



\begin{frame}{Und jetzt?}
    \small 
\textcolor{orange}{
\begin{align*}
        &\sup_{h \in \mathcal{I}(\R^m)} \left|\E[h(\mathbf{Q}(\X))] - \E[h(Z)]\right| \\ & \qquad \qquad \qquad \leq \frac{4}{3} \sup_{h \in \mathcal{I}(\R^m)} \left(\left|\E[h_t(\mathbf{Q}(\X))] - \E[h_t(Z)]\right| + 2 \sqrt{m}\sqrt{t}\right)
        \end{align*}}
    \begin{itemize}
        \item $h_t$ erfüllt die Bedingungen an $\varphi, \varphi', \varphi'', \varphi'''$ für jedes $t$. 
        \item Rechne nach, dass $\norm{\varphi''}_\infty, \norm{\varphi'''}_\infty \leq t^{-\frac{3}{2}}$. 
        \item Seien $B_1$ und $B_2$ die Bounds aus dem letzten Beweis. Dann erhalte
        \begin{align*}
        &\sup_{h \in \mathcal{I}(\R^m)} \left|\E[h(\mathbf{Q}(\X))] - \E[h(Z)]\right| \\ & \qquad \qquad \qquad \leq \frac{8}{3}\sqrt{m}\sqrt{t} + \frac{4}{3}\left(B_1 + B_2\right)t^{-\frac{3}{2}}
        \end{align*}
        \item Ableiten (zweimal) ergibt Minimum für $t = \sqrt{3(B_1 + B_2)/(2\sqrt{m})}$. 
    \end{itemize}
\end{frame}

\begin{frame}{Und jetzt?}
    \small 
    \begin{itemize}
        \item Erhalten
        \begin{align*}
        &\sup_{h \in \mathcal{I}(\R^m)} \left|\E[h(\mathbf{Q}(\X))] - \E[h(Z)]\right| \\ & 
        \leq \frac{8}{3}\sqrt{m}\sqrt{\sqrt{3(B_1 + B_2)/(2\sqrt{m})}} + \frac{4}{3}\left(B_1 + B_2\right)\left(\sqrt{3(B_1 + B_2)/(2\sqrt{m})}\right)^{-\frac{3}{2}} \\ & \sim (B_1+B_2)^{\frac{1}{4}} \to 0.
        \end{align*}
        \item Gilt auch für andere Kovarianzmatrizen.
        \item Damit gilt insbesondere $d_\text{Kol}(\mathbf{Q}(\X), Z) \to 0$. 
    \end{itemize}
\end{frame}

%%%%Drumherum

\section{Ausklang}
\begin{frame}{Was war geschehen?}
    \begin{itemize}
        \item Autoren haben mit $\mathrm{Inf}_i(f)$ gearbeitet. Tauchen auf, wenn man sukzessive $(G_1, \ldots, G_N)$ zu $(X_1, \ldots, X_N)$ transformiert. \pause \footnotesize
        \begin{align*}\E[(V_j^{(i)})^2] &= \sum_{\substack{1 \leq i_1, \ldots, i_{d_j} \leq N_j \\ \exists k: i_k = i}} f_j(i_1, \ldots, i_{d_j})^2 \E[(Z_{i_1}^{(i)})^2] \cdots \textcolor{orange}{\widehat{1}_{(i)}} \cdots \E[(Z_{i_d}^{(i)})^2] \\ &= \sum_{\substack{1 \leq i_1, \ldots, i_{d_j} \leq N_j \\ \exists k: i_k = i}} f_j(i_1, \ldots, i_{d_j})^2 = d_j! \mathrm{Inf}_i(f_j) \end{align*} 
        \normalsize
        \item Einflussindizes kontrollieren Abstand von verrauschter und gaußscher Summe.
        \pause
        \item \textit{Treibstoff}: Hyperkontraktivität.
        \pause
        \item Die sind beschränkt durch Kontraktionen, welche man vor allem durch eine Produktformel für mehrfache \textbf{Integrale} erhält. Unsere Summen sind aber mehrfache Integrale. 
        \item Verbindung von verschwindenden Kontraktionen und Konvergenz in Verteilung gegen Normalverteilung bekannt / plausibel.
        
    \end{itemize}
\end{frame}

% \begin{frame}{Was geht noch? Was geht nicht?}
%     \begin{itemize}
%     \item Man kann auch etwas in der 1-Wasserstein-Metrik zeigen (andere Klasse $\H$ von Lipschitz-Testfunktionen)
%     \pause
%     \item Man kann den Abstand zu $\chi^2$-Verteilungen mit ähnlichen Argumenten und Invarianten beschränken. 
%     \pause
%     \item Man kann die Resultate \textbf{nicht} für den Fall $d=1$ nutzen. Grund: $\sum_{i=1}^n \frac{1}{\sqrt{n}} X_i \Longrightarrow \mathcal{N}(0,1)$, aber \[\norm{h \otimes_1 h}_2 = \sum_{i=1}^n (\sqrt{n}^{-1})^2 = 1 \not \longrightarrow 0.\]
%     \end{itemize}
% \end{frame}